\section{Intrinsic local asymptotic minimaxity}
\label{sec:local-minimax}
Fix a datum on the regular quotient of
Theorem~\ref{thm:quotient-chart}. At clock $T_j$ take $n_j$ independent
categorical observations, independently between clocks, where
$N=\sum_j n_j$ and $n_j/N\to w_j>0$. The design weights are known and
$\sum_jw_j=1$. All cell probabilities are positive on this stratum.
Define, on the product simplex tangent space,
\begin{equation}\label{eq:fisher-object}
 H_{P,w}=\diag\big(w_j/(P_j)_{ab}\big)_{j,a,b},\qquad
 I=J^{\mathsf T}H_{P,w}J,\qquad
 V=GI^{-1}G^{\mathsf T}.
\end{equation}
The matrix $V$ is an observable positive definite form on the locally
labelled root coordinates. Its conjugacy class under component
permutations is intrinsic. Put
\begin{equation}\label{eq:gaussian-risk-object}
 \mathfrak R(P,w)=\E\|Z\|_\infty^2,
                  \qquad Z\sim\mathcal N(0,V).
\end{equation}
Thus the local deterministic condition and the statistical constant
will be two functionals of the \emph{same} observable matrix, rather
than unrelated lower and upper certificates.

For $v\in\R^p$, write $\theta_N(v)=\theta_0+v/\sqrt N$. For fixed
$R<\infty$ these parameters are admissible when $N$ is sufficiently
large. Define the local minimax risk
\begin{equation}\label{eq:local-risk-definition}
 \mathfrak R_N(R)=\inf_{\widehat\Lambda_N}
   \sup_{v^{\mathsf T}Iv\le R^2}
   N\E_{\theta_N(v)}
       d_\infty(\widehat\Lambda_N,\Lambda(\theta_N(v)))^2.
\end{equation}
The infimum is over all measurable multiset estimators. An estimator
in this local problem may depend on the centre $\theta_0$ and the
neighbourhood radius, but not on the unknown local parameter $v$.

\begin{theorem}[Exact local statistical constant]\label{thm:local-minimax}
At every datum in the regular quotient, including cross-component
root collisions,
\begin{equation}\label{eq:local-minimax-equality}
 \lim_{R\to\infty}\liminf_{N\to\infty}\mathfrak R_N(R)
 =\lim_{R\to\infty}\limsup_{N\to\infty}\mathfrak R_N(R)
 =\mathfrak R(P,w).
\end{equation}
A locally efficient minimum-distance estimator followed by component
root extraction and sorting attains the upper bound. Moreover,
\begin{equation}\label{eq:metric-risk-comparison}
 \mathfrak c_{H_{P,w}}(P)^2
 \le \mathfrak R(P,w)
 \le 2\{\log(2kd)+1\}\mathfrak c_{H_{P,w}}(P)^2.
\end{equation}
For fixed dimensions the same observable local condition therefore
controls both sides of local asymptotic minimax risk. Neither equality
asserts a risk lower bound for every estimator at the single point
$\theta_0$.
\end{theorem}

The non-differentiability of sorting at a collision requires a separate
argument. We give it before passing to the categorical experiment.
Let a partition of $\{1,\ldots,q_0\}$ define a group $\Pi$ of
permutations within its blocks, and write
\[
 d_\Pi([u],[v])=\min_{\pi\in\Pi}\|u-\pi v\|_\infty.
\]
A quotient point can equivalently be represented by sorting within each
block. The blocks here will be the groups of equal roots at the base
datum, not the latent components.

\begin{lemma}[A Gaussian quotient experiment]\label{lem:gaussian-quotient}
Let $I$ be positive definite and let $G:\R^p\to\R^{q_0}$ be
surjective. In the experiment $X\sim\mathcal N(v,I^{-1})$, $v\in\R^p$,
with target $[Gv]\in\R^{q_0}/\Pi$, one has
\begin{equation}\label{eq:gaussian-quotient-risk}
 \inf_{\widehat t}\sup_{v\in\R^p}
      \E_v d_\Pi(\widehat t,[Gv])^2
       =\E\|\mathcal N(0,GI^{-1}G^{\mathsf T})\|_\infty^2.
\end{equation}
The same constant is approached from below by Bayes risks of finite
priors supported in bounded subsets of $\R^p$.
\end{lemma}
\begin{proof}
The estimator $[GX]$ has loss at most $\|G(X-v)\|_\infty^2$, giving
the upper bound. First consider the labelled linear target $Gv$.
For the prior $v\sim\mathcal N(0,t^2I^{-1})$, the posterior covariance
is $t^2(1+t^2)^{-1}I^{-1}$. A centrally symmetric distribution and a
convex even loss have a risk-minimizing centre: average the losses at
a displacement and its negative and use convexity. The posterior mean
is consequently Bayes for squared sup-norm loss, and its Bayes risk is
\[
 \frac{t^2}{1+t^2}\E\|\mathcal N(0,GI^{-1}G^{\mathsf T})\|_\infty^2.
\]
Letting $t\to\infty$ proves the labelled minimax lower bound. For each
fixed $t$, truncate this prior to growing cubes, and then approximate
it by finitely supported priors. Gaussian likelihoods are continuous in
$v$ in total variation; compactly supported targets allow actions to be
clipped to a common coordinate box without increasing loss. The Bayes
risks of the finite approximations therefore converge to the truncated
Bayes risk. Removing truncation can be checked directly. Clip an estimator for the
truncated prior to the coordinate box containing $G[-r,r]^p$, whose
radius is at most $C_G r$. Its contribution on the omitted prior tail
is at most
\[
 2C_G^2r^2\Pr(\|v\|_\infty>r)
       +2\E[\|Gv\|_\infty^2\mathbf1_{\{\|v\|_\infty>r\}}],
\]
which tends to zero for the fixed Gaussian prior. The Bayes risk of the
untruncated prior is at most the truncated Bayes risk times its retained
mass plus this tail term. This proves the needed lower approximation. Thus finite priors on growing cubes
approach the labelled constant.

It remains to show that the finite permutation quotient cannot reduce
this worst-case constant. Fix a cube $C=[-r,r]^p$ supporting one of
these finite priors. Since $G$ is onto, there is $v_*$ such that the
coordinates of $Gv_*$ are pairwise distinct within every block. For
large enough $M$, the coordinate ranges
\[
 \{(Gv)_a:\ v\in Mv_*+C\},\qquad a\text{ in a fixed block},
\]
are disjoint intervals in a fixed order. Translate the finite prior by
$Mv_*$. On its support, sorting within each block is a fixed permutation
of the labelled target. Given any quotient estimate, sort each block
and clip its successive coordinates to these disjoint intervals. This
cannot increase the error from any ordered target in the support. The
resulting coordinates have a unique labelling, and their quotient loss
is exactly their labelled sup-norm loss. Hence the quotient Bayes risk
on this translated prior is at least the labelled Bayes risk on the
original prior. Translation of a Gaussian location experiment does not
change that latter risk. Allow $r$ and the prior scale to grow. The
translated supports are bounded at each stage, proving both assertions.
\end{proof}

\begin{proof}[Proof of Theorem~\ref{thm:local-minimax}]
\emph{The limit experiment.} Let $e$ index the $m_1m_2$ cells. Write
$J_{je}$ for the row derivative of their probability at clock $j$.
The score sum, normalized by $\sqrt N$, is
\[
 \Delta_N=\frac1{\sqrt N}\sum_{j=1}^\ell\sum_{i=1}^{n_j}
                         \frac{J_{j,X_{ji}}^{\mathsf T}}{P_{j,X_{ji}}}.
\]
It has mean zero, since $\sum_eJ_{je}=0$, and covariance tending to $I$.
Positive cell floors and the analytic chart make all derivatives
bounded on a fixed neighbourhood. Taylor expansion of each log
probability, the law of large numbers and the finite-dimensional
central limit theorem give, uniformly for bounded $v$,
\begin{equation}\label{eq:categorical-lan}
 \log\frac{d\mathbb P_{\theta_N(v)}}{d\mathbb P_{\theta_0}}
       =v^{\mathsf T}\Delta_N-\tfrac12v^{\mathsf T}Iv+o_{\mathbb P}(1),
 \qquad \Delta_N\ \Longrightarrow\ \mathcal N(0,I).
\end{equation}
For completeness, the sum of the third-order log-probability
remainders is $O(N^{-1/2}\|v\|^3)$, and replacing the empirical second
derivative by its expectation costs $o_{\mathbb P}(1)$ on each bounded
set. Differentiating $\sum_ep_{je}=1$ identifies its expected negative
Hessian with $\sum_eJ_{je}^{\mathsf T}J_{je}/p_{je}$. This is precisely
\eqref{eq:fisher-object}. These are the usual local asymptotic normality
calculations \cite{Hajek72,vdV98}, here in the specified observation
geometry.

\emph{Upper bound and moments.} On a small closed coordinate ball about
$\theta_0$, minimize $\|\widehat P-F(\theta)\|_{H_{P,w}}^2$.
For $\widehat P$ sufficiently close to $P$, its Hessian is uniformly
positive definite on a suitably chosen ball; its unique minimizer is
interior and depends smoothly on $\widehat P$. Define the estimator
arbitrarily in $\cK_d$ on the complementary event. The implicit
normal equations yield, uniformly for bounded $v$,
\begin{equation}\label{eq:efficient-expansion}
 \sqrt N(\widehat\theta-\theta_N(v))
   =I^{-1}J^{\mathsf T}H_{P,w}\sqrt N
                       (\widehat P-F(\theta_N(v)))+o_{L^2}(1).
\end{equation}
To justify the $L^2$ assertion, cell-frequency fourth moments are
$O(N^{-2})$ and the local Taylor remainder is quadratic. The probability
of leaving the fixed neighbourhood is exponentially small, uniformly
on each bounded local parameter set; bounded targets control the
fallback contribution after multiplication by $N$.
The block covariance of $\sqrt N\widehat P$ tends to
\[
 \Sigma_w=\bigoplus_jw_j^{-1}
                    \{\diag(P_j)-\operatorname{vec}P_j\operatorname{vec}P_j^{\mathsf T}\}.
\]
On the simplex tangent, $J^{\mathsf T}H_{P,w}\Sigma_wH_{P,w}J=I$.
Thus the labelled root error obtained from \eqref{eq:efficient-expansion}
converges, with second moments, to $\mathcal N(0,V)$. Matching the
labelled roots and then sorting cannot increase its sup norm, proving
the upper bound in \eqref{eq:local-minimax-equality} for every fixed
$R$. The construction is a local statistical estimator, not a claim of
global optimization complexity. A consistent minimum-distance selector
and localization in the recovered quotient chart give the same
expansion without supplying a latent labelling.

\emph{Lower bound at collisions.} Partition the base aggregate roots
into equal-value groups and let $\Pi$ be their within-group permutation
group. Taylor expansion gives local root velocities $Gv$. For every
fixed finite set of local parameters, their scaled bottleneck losses
converge uniformly, on bounded local action sets, to
$d_\Pi([u],[Gv])^2$. Indeed the different base-value groups are separated
by a fixed gap. Sort any estimated multiset and clip the coordinates
in each base-value group to a common interval of radius $C/\sqrt N$
about that base value, containing all targets of the finite experiment.
Coordinatewise clipping cannot increase the real ordered-matching loss.
After centering and scaling this restricts actions to a compact set;
within each group the limiting operation is exactly the quotient by
$\Pi$, not differentiation of the sorting map.

For a finite prior on local parameters, \eqref{eq:categorical-lan}
implies joint convergence of its likelihood ratios to those of the
Gaussian shift with information $I$. The likelihood ratios are
uniformly integrable: their nonnegative limits have expectation one,
as do the ratios themselves. On the compact action set just described,
losses converge uniformly and are bounded. The formula for posterior
Bayes risk as the integral of the minimum of weighted losses therefore
gives convergence of the finite-experiment Bayes risks. This is the
finite-prior form of the local asymptotic minimax argument, and does
not require an exchange of the limits $N$ and $R$.

By \eqref{eq:root-coefficient-derivative}, $G$ is onto, so
Lemma~\ref{lem:gaussian-quotient} supplies finite priors with Gaussian
Bayes risks arbitrarily close to \eqref{eq:gaussian-risk-object}. Each
prior is contained in an information ball of some finite radius $R$.
Use that prior first, let $N\to\infty$, and only then let $R\to\infty$.
A minimax risk dominates every such Bayes risk. This proves the lower
bound in \eqref{eq:local-minimax-equality}, including at cross-component
collisions.

Finally let $\sigma^2=\max_aV_{aa}=\mathfrak c_{H_{P,w}}(P)^2$.
One coordinate gives $\E\|Z\|_\infty^2\ge\sigma^2$. The union bound
$\Pr(\|Z\|_\infty>t)\le2q_0\exp(-t^2/(2\sigma^2))$, integrated
above $t^2=2\sigma^2\log(2q_0)$, gives the upper bound in
\eqref{eq:metric-risk-comparison} without any independence assumption
on the coordinates of $Z$.
\end{proof}

The theorem supplies intrinsic sharpness on the full open identifiable
quotient stratum, not merely on a specially chosen testing family.
It complements, and does not replace, the closed-model confidence sets,
signal-class lower bounds, and internal-multiplicity neighbourhood rates
below. It uses a fixed limiting allocation $w$; no optimal adaptive
allocation constant is being inferred from it. The adaptive lower
bounds in the later sections retain their stated scope.
